% This version of CVPR template is provided by Ming-Ming Cheng.
% Please leave an issue if you found a bug:
% https://github.com/MCG-NKU/CVPR_Template.

% \documentclass[review]{cvpr}
\documentclass[final]{cvpr}

\usepackage{times}
\usepackage{epsfig}
\usepackage{graphicx}
\usepackage{amsmath}
\usepackage{amssymb}
\usepackage{float}
\usepackage{overpic}
\usepackage{caption}

% Include other packages here, before hyperref.

% If you comment hyperref and then uncomment it, you should delete
% egpaper.aux before re-running latex.  (Or just hit 'q' on the first latex
% run, let it finish, and you should be clear).
\usepackage[pagebackref=true,breaklinks=true,colorlinks,bookmarks=false]{hyperref}

% \def\cvprPaperID{****} % *** Enter the CVPR Paper ID here
% \def\confYear{CVPR 2021}
%\setcounter{page}{4321} % For final version only

\begin{document}

%%%%%%%%% TITLE
\title{Short Title: with a long title describing the key details of this work.}

\author{Student 1\\
% For a paper whose authors are all at the same institution,
% omit the following lines up until the closing ``}''.
% Additional authors and addresses can be added with ``\and'',
% just like the second author.
% To save space, use either the email address or home page, not both
MS in Information Systems\\
Northeastern University\\
student@northeastern.edu
\and{Student 2}\\
MS in Information Systems - Bridge\\
Northeastern University\\
student@northeastern.edu
}

% \maketitle
\twocolumn[{
    \renewcommand\twocolumn[1][]{#1}%
    \maketitle
    \begin{center}
      \centering
      \begin{overpic}[width=\textwidth]{fig/cover_img.jpg}
      \end{overpic}
      \captionof{figure}{Add a figure title and description, and include a reference if the image is sourced online for illustration purposes. This cover page aims to showcase the graphics of your project and capture the readers' attention.} 
      \label{fig_teaser}
    \end{center}
}]


%%%%%%%%% BODY TEXT
\section{Abstract}
Write a 100-word summary providing an overview of this project. Please include: (1) the background and problem narrative, (2) the new system you propose to address this problem, (3) the overall inputs and outputs, and (4) the significance and real-world application scenarios of your work.

%-------------------------------------------------------------------------
\section{Background}
This section aims to illustrate the context of your project. It should include at least two subsections: one is the target market, which explains the industry context, market needs, and existing challenges, and another describing the target users, outlining who will benefit from your solution and what problems or needs they face. You may also add additional subsections, such as technology trends, existing solutions, or research gaps, to help further define the scope and focus of your work. The following sections in this document are provided for illustration purposes only. Please use citations to mention your data sources (\cite{zillow2020home,Curbio2023What,Barbara2022Will})
\\

\textbf{Subsection}  Hui \etal\cite{hui2017illuminant} demonstrated that a illuminant separation problem can be solved if one resolves the reflectance chromaticity of all scene points, which can be calculated by taking a second flash image of the scene. They showed in a later study \cite{hui2019learning} that they could computationally estimate, via a deep neural network, the reflectance chromaticity using one image only. However, Hui solved the problem for separating two illuminants only, in which case it is not required for the neural network to be able to model the latent spectrum distributions of multiple illuminants.

% \maketitle
\twocolumn[{
    \renewcommand\twocolumn[1][]{#1}%
    \maketitle
    \begin{center}
      \centering
      \begin{overpic}[width=\textwidth]{fig/system_diagram.png}
      \end{overpic}
      \captionof{figure}{The system diagram illustrates the system’s inputs and outputs, the internal components that comprise the system, and the interconnections among those components (represented by arrows and lines).}
      \label{fig_teaser}
    \end{center}
}]

\textbf{Potential Users } Prior to Hui \etal\cite{hui2019learning}, instead of regressing the single-illuminant images, a number of methods have been demonstrated to parametrically model the illuminants, such as Lalonde’s method \cite{holdgeoffroy2018deep} utilized Convolutional Neural Network to extract High Dynamic Range (HDR) panoramic images from Low Dynamic Range (LDR) images and predict the virtual objects’ appearance in the image. Gardner \etal\cite{gardner2017learning} proposed a deep neural network to predict lighting conditions under High Dynamic Range (HDR) images, from a single Low Dynamic Range (LDR) image. Their method can recover invisible light sources in the photo without extra scene information.\\


%-------------------------------------------------------------------------



\section{Motivation and Technical Contribution}
\label{sec:motivation}
\textbf{Motivation} The previous research by Hui \etal\cite{hui2019learning} directly regressed the image into two single-illuminant images. As a step forward, we aim to be able to model different spectra distributions with our network, so that we can regress the input image into multiple single-illuminant images. These illuminants, although having entirely different spectra profiles, have wavelength overlaps. We believe this approach can work because (1) this is still an inverse-rendering problem with a step forward from Hui \etal\cite{hui2019learning}, (2) neural networks such as variational autoencoder have been shown to be able to model latent distributions.\\

\textbf{Technical Contribution} Real world scenes are lit by multiple illuminants having different spectral profiles. Isolating those illuminants would require large computational efforts tailored against each individual scene. Our method allows an efficient decrease in computation cost compared to traditional methods, and an effective aid in photometric analysis, as well as interior lighting design.

%-------------------------------------------------------------------------
\section{Challenges}
\label{sec:challenges}

Learning outcomes of the project includes, but are not limited to: (1) large-scale physics-based rendering with Mitsuba in XML; (2) designing the image formation model; (3) translating the image formation model into the sequence of inference steps for learning-based neural networks that can model the latent spectral distributions of lights; (4) designing the evaluation metrics. To achieve these learning outcomes, we will need to learn XML, Mitsuba, and physics-based rendering theories and frameworks.


\section{Project Timeline}
\begin{itemize}
    \item By Feb 22st: finalize research questions and preliminary workflow
    \item By Mar 10th: finish literature review (1-page related work section)
    \item By Mar 24th: generate the preliminary training dataset
    \item By Mar 24th: generate the preliminary training dataset
    \item By Mar 31th: derive our own image formation model and build our deep neural network model (2-page methods section)
    \item By  Apr 7th: finish rendering the entire dataset and get preliminary results
    \item By Apr 21th: test the results with the real-world images (as bonus)
    \item By May 7th: project presentation, full paper completion
\end{itemize}

\section{Result Evaluation}
We will test our framework by using an untouched portion of the synthetic dataset. If time allows, we will capture real-world images and test on that as well. For evaluation purposes, we will provide visual comparison analysis as well as accuracy metrics comparison for the test result against the ground truth.

\section{Expected Outcome}

\textbf{Minimum Goal:} isolate multispectral illuminants from a single image using an existing but small dataset from Smagina et al (24 scenes, 18 color lights)

\textbf{Additional Goal 1:} render our own large-scale dataset by physically modeling illuminants of different spectra

\textbf{Additional Goal 2:} have our method working on our large-scale dataset (1000+ scenes, different kinds of light)

\textbf{Additional Goal 3:} test our method on real-world data


\section{Author Contributions}
\textcolor{blue}{Please include this section if your work as a team}.
Conceptualization, G.J.,J.T.; methodology, G.J.,J.T.; software, G.J.,J.T.; validation, G.J.,J.T.; formal analysis, G.J.,J.T.; writing, review and editing, G.J.,J.T.; visualization, G.J.,J.T.; All authors have read and agreed to the published version of the manuscript.



%-------------------------------------------------------------------------

{\small
\bibliographystyle{ieee_fullname}
\bibliography{egbib}
}

\end{document}
